% MedVision manuscript -- Introduction section

\section{Introduction}\label{sec:introduction}

Medical vision--language models (VLMs) combine radiographs with the clinical reports
that accompany them. This pairing is valuable because images encode visual
structure whereas reports provide findings and context; resources such as
MIMIC-CXR make both modalities available at scale \citep{johnson2019mimiccxr}.
High multimodal accuracy, however, does not establish that the model used
clinically meaningful visual evidence. A model may exploit a report shortcut or
another predictive feature that is easier to optimise than image information.

This distinction is especially important when targets are report-derived. The
CheXpert labeler converts report language into uncertainty-aware
chest-radiograph labels \citep{irvin2019chexpert}. If the same report, or its
\texttt{findings} field, is also supplied to the model, the target and one
input modality share an information source. The resulting task may reward
linguistic recognition of the report-derived label rather than independent
image-grounded recognition. Report-conditioned prediction can still be useful;
the problem is claiming incremental image evidence from aggregate accuracy
alone.

Recent evidence shows why this distinction matters in medical VLM evaluation.
Performance can change with the amount of clinical context supplied before an
image is seen, and apparently correct multimodal answers may be accompanied by
unreliable visual rationales \citep{wang2026clinicalcontext,jin2024hiddenflaws}.
Large-scale chest-radiography audits have likewise found that removing the
image can leave performance largely intact, while consistency across prompts
can coexist with non-visual decision making \citep{lotfinia2026needimage,sadanandan2026consistentdangerous}.
Counterfactual prototype methods can improve interpretability, but they do not
replace a report-preserving intervention that directly measures the change in
the prediction \citep{chiang2026cxrnexus}.

Conventional model comparisons rarely replace the image while holding the
report, label, architecture, and evaluation protocol fixed. Similarly, a
non-zero fusion gate or attention map is not necessarily causal evidence of
visual use: gates can remain near initialization and a one-token interaction
cannot show word--region alignment. Text-debiasing methods may therefore lower
a score without creating an independent visual signal when the endpoint itself
is linked to text.

We present a controlled audit of a three-class MIMIC-CXR-derived task
(Normal, Pneumonia, and Pneumothorax). The model receives a frontal radiograph
and the report's \texttt{findings} field. We compare full multimodal models
with a report-preserving Gray-Square control that replaces every radiograph by
a constant tensor, together with text-only and image-only controls. Four
fusion or text-debiasing variants are evaluated alongside the project's
baseline gated dual-path model (CADQ). We
report discrimination, calibration, class-specific behaviour, learned image
gates, and qualitative gradient-weighted class activation mapping (Grad-CAM)
checks. The primary diagnostic is the Image
Contribution Metric (ICM), defined as the relative Macro-F1 difference between
the full model and the Gray-Square reference; paired inference is performed at
the trained-seed level because held-out studies recur across seeds.

The audit asks whether removing image information changes performance, whether
the interventions increase incremental image contribution, whether checkpoint
gates show learned reliance or stagnation, and whether calibration and
class-specific results qualify the headline metrics. It contributes a
report-preserving image ablation, seed-level uncertainty and calibration
reporting, and an explicit analysis of label provenance and control asymmetry.
The study is a diagnostic evaluation, not a new clinical predictor or a claim
that radiographs are generally unnecessary. Its interpretation is limited to
this report-linked cohort, preprocessing path, and checkpoint archive; broader
claims require corrected implementations, independent labels, matched controls,
and external validation. Reporting follows the transparency principles of the
Checklist for Artificial Intelligence in Medical Imaging (CLAIM) and the
TRIPOD+AI framework \citep{mongan2024claim,collins2024tripodai}.
